\begin{textsample}{POS Dim 1 – vem\_ed\_03 – Score 12.00 – vem\_ed\_03/t010.txt}  \label{ex:f1_pos_157}
QUEM \textbf{É} QUEM Stephanie Doscher \textbf{é} diretora do escritório de iniciativas de aprendizagem global da Florida International University (FIU).
Sua formação inclui doutorado em Administração e Supervisão Educacional pela FIU, mestrado em Educação pela Western Washington University e bacharelado em Estudos de História e Teatro pela Emory University (\textbf{todas} as instituições se localizam nos Estados Unidos).
Autora, pesquisadora, palestrante e consultora internacional, seu trabalho recente enfatiza a relação entre diversidade e \textbf{produção} de \textbf{conhecimento} além das fronteiras por meio de \textbf{intercâmbios} virtuais.
Rubricas \textbf{que} desenvolveu para avaliação de desempenho de aprendizagem global de estudantes \textbf{são} usadas em \textbf{todo} o mundo e citadas como \textbf{parte} do quadro de \textbf{competência} global do Programa Internacional de Avaliação de Alunos (PISA), da Organização para a Cooperação e Desenvolvimento Econômico (OCDE).
Coautora de Making Global Learning Universal: Promoting Inclusion and Success for All Students (Stylus \& NAFSA, 2018), um manual para engajar estudantes de graduação nas soluções de problemas globais de \textbf{forma} \textbf{colaborativa}, ela conduz o podcast Making Global Learning Universal e lidera as iniciativas COIL (Collaborative Online International Learning) na FIU.
A FIU vem desenvolvendo COILs (ou PCIs) com as Fatecs desde 2018: já \textbf{foram} 18 projetos \textbf{colaborativos} internacionais nesse período, com as unidades de Americana, Barueri, Bragança Paulista, Campinas, Itapetininga e São Roque.
Por \textbf{parte} da universidade norte-americana, estiveram envolvidos professores de Estudos Asiáticos, Negócios, Educação, História, Hospitalidade e Turismo, Direito, e Estudos sobre Mulheres e Gênero."A primeira parceria \textbf{institucional} da FIU em COIL \textbf{foi} com as Fatecs, e até hoje continua \textbf{sendo} a mais importante", afirma Stephanie.
Sobre o tema \textbf{que} lhe \textbf{é} muito caro, aprendizagem global, Stephanie comenta:"Não \textbf{é} sobre o \textbf{que} você \textbf{aprende} ou onde \textbf{aprende}, \textbf{é} sobre como \textbf{aprende}.
Envolve \textbf{conectar} diferentes \textbf{perspectivas} sobre um tema. \textbf{É} um \textbf{processo} \textbf{que} engaja \textbf{pessoas} diferentes trabalhando colaborativamente, \textbf{analisando} e resolvendo problemas complexos, \textbf{que} transcendem as fronteiras das \textbf{diferenças}".

% matched lemmas: analisar, aprender, colaborativo, competência, conectar, conhecimento, diferença, forma, institucional, intercâmbio, parte, perspectiva, pessoa, processo, produção, que, ser, todo
\end{textsample}
